%& -shell-escape
% !TeX root = thesis-text.tex
% !TeX encoding = UTF-8
% !TeX program = XeLaTeX

\documentclass{.local/lib/classes/ndvr-super-article-standalone}

\begin{document}

    \section{Постановка проблемы}

    % Связанные с~видео онлайн-услуги,
    % такие как~публикация и~распространения видео,
    % трансляция видео, видео-рекомендации и~пр,
    % все более и~более смещают сферу интересов пользователей.
    % Сейчас пользователи не~только могут искать и~просматривать видео,
    % но~и~заниматься редактированием, размещением и~вещанием.

    \subsection{Рост потребления видео}

    Согласно данным аналитической компании «ComScore»
    \index{comScore}\footnote{
        ComScore (ориг. «comScore») – 
        агентство, занимающееся анализом рынка интернет-технологий.
        Официальный сайт: \href{http://comScore.com}{comscore.com}.
    },
    Россия — ведущая страна Европы по~зрителям интернет-видео
    и~вторая по~числу просмотров в~месяц \cite{Henni:2013}.

    Количество видео размещенных на~русскоязычных
    сайтах за~период с~июня 2012 до~июня 2013 года выросло
    на~16\%, с~7,9 миллионов до~9,1 миллионов.

    В 2013 году число российских интернет-зрителей 
    старше 15 лет превысило 55~миллионов\cite{Onishenko:2013}. 
    В январе 2017 года общее число уникальных 
    посетителей видео-ресурсов превысило 74,4~миллиона\footnote{
        \href{https://adindex.ru/news/researches/2017/03/16/158635.phtml}
        {https://adindex.ru/news/researches/2017/03/16/158635.phtml}  
    }.

    По~данным компании «Медиаскоп»\index{TNS}\footnote{
        Медиаскоп~— международная исследовательская группа,
        занимающаяся изучением аудитории СМИ в России.
        Официальный сайт: \href{http://mediascope.net}{mediascope.net}.
    },
    просмотр видео~— второй по~популярности
    вид деятельности после социальных сетей \cite{Henni:2013}.
    В~2012, он занимал около~10\% от~общего 
    времени проведенного в~Интернете. 
    В январе 2017 года~— уже 30\%.


    По данным компании Cisco\footnote{
        Cisco—американская транснациональная компания, разрабатывающая и продающая сетевое оборудование, предназначенное в основном для крупных организаций и телекоммуникационных предприятий.
    }, к 2019 году интернет-видео составит 80\% 
    от мирового интернет-трафика\footnote{
        Cisco Visual Networking Index: 
        Forecast and Methodology, 2016–2021.
        Updated:September 15, 2017
        Document ID:1465272001663118
        https://www.cisco.com/c/en/us/solutions/collateral/service-provider/visual-networking-index-vni/complete-white-paper-c11-481360.html
    }.

    \begin{figure}[H]
        \includegraphics[width=\textwidth]
            {.local/vec/out/pdf/video-marker-future.pdf}
        \caption{Ретроспектива и~прогноз 
            развития российского рынка онлайн-видео
            в~2012—2018 годах. Данные взяты из~работы 
            \cite{Henni:2013}
        }
        \label{fig:video-marker-future}
    \end{figure}


    Обычные СМИ осваивают Интернет и~предоставляют свою продукцию
    потребителям веб-приложений и~сервисов.
    Музыкальные клипы публикуются 
    на~официальных каналах в~YouTube\footnote{
        YouTube — видеохостинг, предоставляющий 
        услуги хранения, доставки и показа видео. 
        Официальный сайт: \href{https://www.youtube.com/}{youtube.com}.
    }.
    Телевизионные программы транслируются на~сайтах телекомпаний.
    Трейлеры фильмов, в~том числе еще~не~вышедших в~прокат,
    предварительное размещаются  на~тематических сайтах.

    Быстрый рост видеоприложений
    и~услуг впоследствии приводит к~росту онлайн-видео,
    как~отмечено в~работах \cite{Tan:2009}, \cite{Shen:2005}, \cite{Wu:2007a}.


    \subsection{Рост количества нечётких дубликатов}

    Издатели, пользователи и~сторонние организации часто редактируют
    и~публикуют копии оригинального видео.
    Сходные видео, на~которых изображено одно и~то~же событие,
    могут быть приобретены и~опубликованы многократно.

    Все эти явления способствуют существенному росту процента
    нечетких дубликатов видео в~Интернете.
    Количество дубликатов достигает 93\% выдачи
    по~некоторым поисковым запросов \cite{Wu:2007a}.

    Большое количество нечётких дубликатов 
    видео предъявляет высокие требования
    к~их~поиску.
%     Поиск нечетких дубликатов видео применяют для
%     \begin{enumerate}


        
    % Это важно для~многих сфер применения,
    % таких как~обнаружение нелицензионных копий,
    % видеонаблюдение, ранжирование 
    %Интернет-видео, видео рекомендация, и~т.~д.

    %В~следующих разделах мы~опишем ряд сценариев подробнее
    %и~покажем как~поиск НДВ формирует базу для~некоторых областей,
    %упомянутых ранее.

\end{document}

